% !PRINTED 345
\weberpage{345}

%<<CONT
in which $\varphi, \psi$ are integral functions of the variables $\xi, \eta$
with integral coefficients, whose degree with respect to neither of the variables
exceeds the value 2. If we order these functions according to $\eta$, we may
set:
\weq{5}{\xi' = A\eta^2 + B\eta + C, \quad \eta' = A'\eta^2 + B'\eta + C',}
in which $A, B, C, A', B', C'$ are integral functions of $\xi$, at most
of degree $2$. In the equations (4) we shall hold the
variable $\xi$ fixed and set $\eta = 0, 1, 2$; that is, we apply
(5) to the triple $(\xi, 0)\,(\xi, 1)\,(\xi, 2)$. The corresponding
$(\xi_1', \eta_1'), (\xi_2', \eta_2'), (\xi_3', \eta_3')$ must then satisfy the relations (3),
hence:
\[
\begin{array}{ll}
\xi_1' \equiv C, & \eta_1' \equiv C',\\
\xi_2' \equiv A + B + C, & \eta_2' \equiv A' + B' + C',\\
\xi_3' \equiv A - B + C, & \eta_3' \equiv A' - B' + C',
\end{array}
\]
and from this it follows that $A$ and $A'$ must be congruent to
zero for every $\xi$.

In the same way one may conclude that no term with $\xi^2$ can
occur in the substitutions (4), and that these substitutions
 must therefore have the form:
\[
\begin{array}{l}
\xi' \equiv m\,\xi\,\eta + a\,\xi + b\,\eta + c\\
\eta' \equiv m'\,\xi\,\eta + a'\,\xi + b'\,\eta + c'
\end{array} \quad (\mathrm{mod}\ 3).
\]
But if one applies these substitutions to the triple $(0,\,0)\,(1,\,1)\,(2,\,2)$,
and sets the sum of the corresponding $\xi'$ and $\eta'$ equal to
zero, it follows that we must have $m = m' = 0$, and hence that
every permutation of the group $P$ must be contained in the form
\weq{6}{\begin{array}{l}
\xi' \equiv a\,\xi + b\,\eta + c\\
\eta' \equiv a'\,\xi + b'\,\eta + c'
\end{array} \quad (\mathrm{mod}\ 3)}
as was asserted.

\gesperrt{2. Conversely, by every linear substitution
of the form (6) three quantities $(\xi, \eta)$ of a triple
pass into three $(\xi', \eta')$, which likewise form a
triple.}

For from (6) it follows, for any three pairs of numbers $(\xi_1, \eta_1)$,
$(\xi_2, \eta_2), (\xi_3, \eta_3)$:
\weq{7}{\begin{array}{l}
\xi_1' + \xi_2' + \xi_3' \equiv a\,(\xi_1 + \xi_2 + \xi_3) + b\,(\eta_1 + \eta_2 + \eta_3)\\
\eta_1' + \eta_2' + \eta_3' \equiv a'\,(\xi_1 + \xi_2 + \xi_3) + b'\,(\eta_1 + \eta_2 + \eta_3)
\end{array} \quad (\mathrm{mod}\ 3);}
%CONT>>
